\begin{lstlisting}[language={}]
friedman did some genetics work for him, but, because he was handywith a camera, he helped
the cryptologists who were looking for bacon'scipher signatures in shakespeare by making
photographic enlargementsof the elizabethan printing that figured in the work. the
department ofciphers of the riverbank laboratories consisted of 14 or 15 high-schooland
college graduates who assigned the individual letters in theseelizabethan texts to one or
the other of two fonts of type as part of thebaconian search. fabyan gave them their living
plus a salary of about$50 a month. the staff was fed and housed in engledew and
hoovercottages, the cipher laboratories taking up the first floor of engledew.the young
woman who collated the work of many of the other staffmembers was elizebeth smith. she had
been , born august 26, 1892, inhuntington, indiana, the youngest of the nine children of
john m. smith,a dairyman, banker, and county republican committeeman, and hiswife, sopha,
who spelled her daughter's christian name witb an einstead of an a in the middle because she
was not going to have anyonecalling her child "eliza." after completing high school in
huntington,elizebeth attended wooster college briefly but was graduated fromhillsdale
college in michigan where she had majored in english. whileworking at the newberry library
in chicago, she was recruited byfabyan and began work there in 1916.neither she nor friedman
had given any particular previous thoughtto cryptology, but they began to get personally
interested in the work. itis yet another of the ironies of cryptologic history that the
interest of twoforemost cryptologists was aroused by a false doctrine--a doctrine,moreover,
against which they later were to wage a lifetime battle. for attable at the riverbank
cottages they heard gaudy tales of lustyelizabethan life, of the not-so-virgin queen, of
courtiers' intrigues andthe secret histories of the great names of english history--all
actuallyinvalid decipherments of shakespeare's plays tending to prove thatbacon had written
them, related by the gentle, upright, but self-deludedwoman who had "deciphered" them, mrs.
elizabeth wells gallup. thesestories stirred friedman's dormant interest; he began to do
some of thecryptology, and inevitably its puissant magic seeped like the fume ofpoppies into
his mind and spirit and intoxicated him."when it came to the cryptology," he recalled years
later,"something in me found an outlet."an understatement. he soon found himself head of the
department ofciphers as well as the department of genetics at riverbank. theattraction he
felt for cryptology was reinforced by the attraction he felt fora cryptologist: the quick-
witted and sprightly miss smith. in may of 1917they were married and started the most famous
husband-and-wife teamin the history of cryptology.america had declared war a month before,
and river-bank, which hadthe only going cryptologic concern in the country, began getting,
on aninformal basis, cryptograms for solution from various governmentbureaus. probably the
most important were messages to and from a ringof 125 hindus who, with german aid, were
taking advantage ofengland's preoccupation in europe to strike for indian independence.the
intercepts were given to friedman for solution, and he quickly solvedthe number cipher used
in cablegrams to berlin. the letters of theplaintext and of the keyword were transformed
into digits by means of a4 x 7 checkerboard with a normal alphabet;  the key digits were
thenadded to those of the plaintext to form the ciphertext. one key waslamp. each agent had
his own key, but friedman had no trouble insolving them. nor was he stumped by a system
usually regarded byamateurs as the ne plus ultra of cryptographic security: a book cipher.
\end{lstlisting}
